\chapter{Homeostasis}\label{ch:g9:homeostasis}

Your internal conditions stay within narrow ranges even when the outside
world swings. \omterm{def:g9:homeostasis:def}{Homeostasis} is that stability --- achieved by sensors,
\omterm{def:g9:homeostasis:loop}{control centres} and \omterm{def:g9:homeostasis:loop}{effectors} in feedback loops.

\section{The idea}

\begin{definition}[Homeostasis]\label{def:g9:homeostasis:def}
\emph{Homeostasis}\index{homeostasis} is the maintenance of a stable
internal \omterm{def:g4:traits:env}{environment} (temperature, \omterm{def:g4:plants-need:needs}{water} balance, \omterm{def:g7:digestion:blood}{blood} glucose, pH,
etc.) within limits compatible with life.
\end{definition}

\begin{definition}[Feedback loop parts]\label{def:g9:homeostasis:loop}
A \omterm{def:g6:method:control}{control} system typically has: a \emph{receptor}\index{receptor
(homeostasis)} (detects change), a \emph{control centre}\index{control
centre} (compares to a set point), and an \emph{effector}\index{effector}
(makes the adjustment).
\end{definition}

\begin{omfigure}
\begin{tikzpicture}[font=\small,
  b/.style={draw, rounded corners, align=center, minimum width=2.0cm,
    minimum height=0.7cm}]
  \node[b, fill=omDef!15] (r) at (0,0) {receptor};
  \node[b, fill=omProp!15] (c) at (3.5,0) {control\\centre};
  \node[b, fill=omMeth!15] (e) at (7,0) {effector};
  \draw[-Stealth, thick] (r) -- (c);
  \draw[-Stealth, thick] (c) -- (e);
  \draw[-Stealth, thick] (7,-1.0) -- (0,-1.0) -- (0,-0.4);
  \node at (3.5,-1.4) {response alters condition (feedback)};
\end{tikzpicture}

{\small General negative-feedback architecture.}
\end{omfigure}

\section{Negative feedback}

\begin{definition}[Negative feedback]\label{def:g9:homeostasis:neg}
\emph{Negative feedback}\index{negative feedback} reverses a change:
if a variable rises above the set point, \omterm{def:g7:nervous:stim}{responses} lower it; if it falls,
\omterm{def:g7:nervous:stim}{responses} raise it. Most homeostatic \omterm{def:g6:method:control}{controls} are negative feedback.
\end{definition}

\begin{example}[Body temperature]\label{ex:g9:homeostasis:temp}
Too hot: sweat and increased \omterm{def:g5:nerves:senses}{skin} \omterm{def:g7:digestion:blood}{blood} flow lose heat. Too cold:
shivering generates heat; \omterm{def:g5:nerves:senses}{skin} vessels constrict to reduce loss.
\end{example}

\begin{example}[Blood glucose]\label{ex:g9:homeostasis:glu}
After a meal, insulin promotes glucose uptake/storage, lowering \omterm{def:g7:digestion:blood}{blood}
glucose. When glucose is low, glucagon promotes release from stores ---
opposing hormones in a balance.
\end{example}

\begin{omfigure}
\begin{tikzpicture}[font=\footnotesize]
  \draw[thick, -Stealth] (0,0) -- (6.5,0) node[right] {time};
  \draw[thick, -Stealth] (0,0) -- (0,3) node[above] {variable};
  \draw[dashed, omThm] (0,1.5) -- (6.2,1.5);
  \node[omThm] at (5.3,1.8) {set point};
  \draw[thick, omExo] (0.3,1.5) .. controls (1,2.4) and (1.5,2.4)
    .. (2.2,1.5) .. controls (3,0.6) and (3.5,0.6) .. (4.2,1.5)
    .. controls (5,2.2) and (5.5,2.0) .. (6.0,1.55);
\end{tikzpicture}

{\small \omterm{def:g9:homeostasis:neg}{Negative feedback} pulls the variable back toward the set point
after disturbances.}
\end{omfigure}

\section{Water and waste}

\begin{proposition}[Kidneys and balance]\label{prop:g9:homeostasis:kidney}
Kidneys adjust \omterm{def:g4:plants-need:needs}{water} and salt reabsorption and excrete urea. Thirst and
hormones help match intake and loss to keep \omterm{def:g7:microscope:cell}{cells} from swelling or
shrinking osmotically.
\end{proposition}
\admitted

\begin{method}[Analysing a homeostatic story]\label{met:g9:homeostasis:anal}
\begin{enumerate}
  \item Name the variable.
  \item State the disturbance.
  \item Identify sensor, \omterm{def:g6:method:control}{control} and \omterm{def:g9:homeostasis:loop}{effector} if possible.
  \item Say whether feedback is negative or positive.
\end{enumerate}
\end{method}

\section{Exercises}

\begin{exercise}[$\star$]\label{exo:g9:homeostasis:1}
Define \omterm{def:g9:homeostasis:def}{homeostasis}. Give two variables that are kept stable.
\end{exercise}

\begin{exercise}[$\star$]\label{exo:g9:homeostasis:2}
Name the three parts of a basic \omterm{def:g6:method:control}{control} loop.
\end{exercise}

\begin{exercise}[$\star$]\label{exo:g9:homeostasis:3}
What is \omterm{def:g9:homeostasis:neg}{negative feedback}?
\end{exercise}

\begin{exercise}[$\star$]\label{exo:g9:homeostasis:4}
Describe one \omterm{def:g7:nervous:stim}{response} to overheating.
\end{exercise}

\begin{exercise}[$\star$]\label{exo:g9:homeostasis:5}
How do insulin and glucagon differ in their effects on \omterm{def:g7:digestion:blood}{blood} glucose?
\end{exercise}

\section{Problem: Feedback engineer}

